\plannedchapter{Pipeline、前端和后端总览}
{本章从极简 CPU 过渡到现代乱序 CPU，把 front-end、back-end、retire 和 memory subsystem 放到同一张图里。}
{\item 指令从字节到执行结果经过哪些阶段？
\item 前端和后端分别可能如何成为瓶颈？
\item 为什么 ISA 指令不是 CPU 内部执行的最小单位？}
{\item fetch、decode、uop cache、rename、dispatch、scheduler、execute、retire。
\item pipeline hazard、dependency、flush。
\item retire in order 与 execute out of order。}
{用简单汇编和 \code{llvm-mca} 观察指令吞吐，画出一个 loop 的前后端数据流。}
